\documentclass[11pt]{article}
\usepackage{../shared/hanzocover}
\usepackage[utf8]{inputenc}
\usepackage[T1]{fontenc}
\usepackage{amsmath,amssymb}
\usepackage{graphicx}
\usepackage{booktabs}
\usepackage{xcolor}
\usepackage{listings}
\input{../shared/lstlang}
\usepackage{tikz}
\usetikzlibrary{shapes,arrows.meta,positioning,fit,backgrounds,calc}
\usepackage[hidelinks]{hyperref}

\definecolor{codegreen}{rgb}{0,0.6,0}
\definecolor{codegray}{rgb}{0.5,0.5,0.5}
\definecolor{codepurple}{rgb}{0.58,0,0.82}
\definecolor{backcolour}{rgb}{0.96,0.96,0.94}
\definecolor{hzblue}{rgb}{0.12,0.30,0.55}
\definecolor{hzslate}{rgb}{0.22,0.25,0.30}
\definecolor{hzgreen}{rgb}{0.10,0.45,0.25}
\definecolor{hzred}{rgb}{0.60,0.16,0.16}

\lstdefinestyle{mystyle}{
    backgroundcolor=\color{backcolour},
    commentstyle=\color{codegreen},
    keywordstyle=\color{codepurple},
    numberstyle=\tiny\color{codegray},
    stringstyle=\color{codegreen},
    basicstyle=\ttfamily\footnotesize,
    breakatwhitespace=false,
    breaklines=true,
    captionpos=b,
    keepspaces=true,
    numbers=none,
    showspaces=false,
    showstringspaces=false,
    showtabs=false,
    tabsize=2,
    frame=single,
    rulecolor=\color{codegray}
}
\lstset{style=mystyle}

\title{Rebuilding One App in Two Seconds:\\
Run-Time Plugin Decomposition of a 4{,}347-Package Go Binary}
\author{Zach Kelling\thanks{Corresponding author: research@hanzo.ai} \\
    \textit{Hanzo Industries Inc (Techstars '17)} \\
    Los Angeles, California \\
    \texttt{https://hanzo.ai}}
\date{Sequel to HIP-0106 \quad Revision~1}

\begin{document}

\hanzocoverpage

\twocolumn[
  \begin{@twocolumnfalse}
  \maketitle
  \begin{abstract}
  \noindent
  The Hanzo unified cloud binary (HIP-0106) collapsed a microservices control
  plane into one Go process, eliminating the service mesh and the
  container-per-service tax. It also produced a build with 4{,}347 reachable
  packages, and we report a consequence that the original design did not
  anticipate: \emph{editing one line in one of 107 applications costs the same
  as editing all of them}. We measure a one-app rebuild at 20.23\,s with a
  5.6\,GB peak resident set, against a 19.49\,s no-op rebuild --- the marginal
  cost of the edit itself is under a second, and everything else is the link,
  which is a function of total reachable packages and is therefore paid in full
  regardless of what changed. We describe a decomposition that keeps the
  single-artifact deployment while removing the shared link. An application
  builds as its own binary and is composed into a host at run time over a
  unix-domain socket carrying the ZAP binary protocol; the host embeds those
  binaries with \texttt{go:embed} and still ships as one file. The composition
  primitive is a single function type, \texttt{Service = func(*App) error},
  which a linked-in application and a separately built one both inhabit, so
  where a service runs becomes a deployment decision rather than a code change.
  Measured on the same tree and machine, a one-app rebuild falls to 2.11\,s
  (9.6$\times$) with a 0.63\,GB peak resident set (8.9$\times$), and the host
  does not rebuild at all. We explain why Go's \texttt{-buildmode=plugin} cannot
  serve this role, give the three invariants that make repeated hot reloads flat
  in memory, and report a second and larger finding: a single import in the root
  package placed 1{,}154 packages --- every Kubernetes, AWS, gRPC and container
  SDK in the tree --- beneath all 107 applications, and inverting it cut that
  floor from 1{,}798 packages to 644.
  \end{abstract}
  \vspace{1.5em}
  \end{@twocolumnfalse}
]

\section{The cost that was not modelled}

HIP-0106 argued the unified binary on operational grounds: in-process calls
instead of network round-trips, no sidecars, no per-call mutual TLS, one
artifact per brand. Those hold. What the design did not model is that a
statically linked Go binary has a \emph{shared} link, and the link is the
dominant term in the edit--build--test loop.

Table~\ref{tab:baseline} gives the baseline, measured warm on a 20-core
\texttt{aarch64} host with a 28\,GB Go build cache, using
\texttt{/usr/bin/time} for wall clock and peak resident set. Each figure is the
whole \texttt{go build} invocation, which is what a developer actually waits
for.

\begin{table}[h]
\centering
\begin{tabular}{lrr}
\toprule
\textbf{\texttt{cmd/cloud}, 4{,}347 pkgs} & \textbf{wall} & \textbf{peak RSS} \\
\midrule
no-op rebuild            & 19.49\,s & 5.6\,GB \\
pure relink              & 20.33\,s & 5.6\,GB \\
rebuild one small app    & 20.23\,s & 5.6\,GB \\
\bottomrule
\end{tabular}
\caption{The three are the same number. Compilation of the edited package is
noise; the link is the cost.}
\label{tab:baseline}
\end{table}

Two consequences follow. First, the loop does not get faster by making any
individual application smaller, because no application is the cost. Second, the
5.6\,GB peak is per concurrent build: on a machine with 121\,GB of RAM, a
handful of parallel builds is enough to exhaust it, and we observed exactly that
--- an out-of-memory event traced to concurrent Go links against a 24\,GB
\texttt{tmpfs}.

It is worth being precise about what does \emph{not} help. The repository
already generates 87 per-application \texttt{main} packages, each of which
serves exactly one application. They do not help: each still links every
application and disables the rest at run time via an enable-list, so
\texttt{go list -deps ./cmd/product} returns the same 4{,}347 packages as
\texttt{./cmd/cloud}. A runtime enable-list is not a build boundary.

\section{Why not \texttt{-buildmode=plugin}}

Go ships a plugin mechanism, and it is the obvious candidate. It is unusable
here for three independent reasons, each sufficient on its own.

It requires cgo. Building with \texttt{CGO\_ENABLED=0} fails outright with
\emph{``-buildmode=plugin requires external (cgo) linking, but cgo is not
enabled''}. Our images build \texttt{CGO\_ENABLED=0} to \texttt{scratch};
adopting plugins means adopting a libc base image across the fleet.

It requires byte-identical builds of every shared dependency between host and
plugin. With 1{,}923 modules in the graph, that is a coupling stronger than the
static link it replaces --- the host and plugin must be built together to be
loadable together, which is precisely the property we are trying to remove.

It cannot unload. Go has no \texttt{dlclose}: a loaded plugin's code, globals
and goroutines are unreclaimable for the lifetime of the host. A mechanism whose
purpose is to replace a running application with a new build cannot be built on
a primitive that can never release the old one.

A child process has none of these properties. It releases every byte when it
exits, it needs no build-time agreement with its host, and it works under
\texttt{CGO\_ENABLED=0}.

\section{One type for both}

The composition primitive is a single function type:

\begin{lstlisting}[language=Go]
type Service func(*App) error
\end{lstlisting}

A unit of functionality attaches its own routes, middleware and shutdown hooks.
A constructor that takes dependencies and returns one of these is the same thing
curried, so a composition root only ever sees \texttt{Service}. The existing
mount signature in the unified binary, \texttt{func(app *zip.App, deps Deps)
error}, \emph{is} this type curried on \texttt{deps}, so the 107 applications
required no reshaping.

The load path returns the same type, which is the whole point:

\begin{lstlisting}[language=Go]
//go:embed bin/billing
var billingBin []byte

app.Add(
    billing.New(deps),
    zip.Load("/v1/billing", zip.Plugin{
        Name: "billing", Bin: billingBin}),
    zip.Load("/v1/ml", zip.Plugin{
        Name: "ml", Addr: mlAddr}),
)
\end{lstlisting}

The three arguments are, respectively, an application linked into this binary,
one that ships as its own binary embedded in it, and one already running
elsewhere. They have the same type. Moving a service between those lines is a
deployment decision, not a code change --- and critically, the single-artifact
property that motivated HIP-0106 survives, because \texttt{go:embed} puts the
plugin binaries inside the host binary.

Serving and delegating are symmetric, and read the transport off the address
scheme rather than off a method name:

\begin{lstlisting}[language=Go]
app.Listen("/run/hanzo/billing.sock")  // in the plugin
app.Mount("/v1/billing", sock)         // in the host
\end{lstlisting}

The scheme names the protocol; the address names where it is spoken. A path is
a unix socket, a \texttt{host:port} is TCP, so one wire never needs two schemes.
A plugin is an ordinary application: no SDK, no schema, no code generation. Its
only plugin-specific line is reading the address its host assigned.

\section{Reload without leaking}

Replacing a running application with a new build is the operational payoff, and
it is where a naive implementation leaks. Three invariants hold.

\textbf{Routes register once.} They are registered at load and resolve their
target per request through an atomic pointer. Re-registering on each reload
would grow the route table without bound --- the leak that makes hot reload
untenable in practice.

\textbf{The replacement must be listening before any request moves to it.} The
new process is started and proven to accept on its socket while the old one
keeps serving. A build that fails to start changes nothing and returns an error;
a bad deploy cannot take the route down.

\textbf{The old process is released completely.} Killed, reaped exactly once
(one \texttt{Wait} per process, its result delivered on a channel), its pooled
connections closed, its private directory removed.

Unloading leaves the routes registered and answering 503 rather than removing
them, for the same reason: the route table is never mutated after load, so
repeated load/unload cycles are flat.

These are verified rather than asserted. The test suite builds two genuinely
different binaries from identical source, distinguished by a linker-stamped
version, and asserts \emph{which} one answered across four consecutive swaps, a
deliberately failed swap, and an unload/restore cycle. Race detector clean.

\section{Measurements}

Table~\ref{tab:gains} compares the same application under both regimes, on the
same tree, same machine, same warm cache. The plugin is one real application
from the unified binary, built as its own binary; the host is a binary that
links the web framework and the transport and no application at all.

\begin{table}[h]
\centering
\begin{tabular}{lrrr}
\toprule
& \textbf{pkgs} & \textbf{rebuild} & \textbf{peak RSS} \\
\midrule
host only        &   313 &  0.65\,s & 0.31\,GB \\
one app (plugin) &   646 &  2.11\,s & 0.63\,GB \\
monolith         & 4{,}347 & 20.23\,s & 5.6\,GB \\
\bottomrule
\end{tabular}
\caption{Rebuilding one application. The monolith figure is the cost of editing
\emph{any} application; the plugin figure is the cost of editing \emph{that}
application, and the host does not rebuild at all.}
\label{tab:gains}
\end{table}

The edit--rebuild loop improves 9.6$\times$ in wall clock and 8.9$\times$ in
peak memory. Binary size falls from 497.1\,MB to 43.3\,MB for the plugin and
25.4\,MB for the host.

The memory result matters more than the ratio suggests. At 5.6\,GB per link,
concurrency is bounded at roughly twenty parallel builds on a 121\,GB machine
before thrashing, and in practice far fewer alongside anything else. At
0.63\,GB, the same machine sustains an order of magnitude more --- and plugin
builds are mutually independent, so they genuinely parallelise, whereas the
monolith's link is one serial operation no amount of hardware shortens.

Link time scales with reachable packages rather than with edited code, which the
three rows make legible: 313 packages link in 0.65\,s, 646 in 2.11\,s, 4{,}347
in 20.23\,s.

\section{The dependency floor}

Decomposition exposed a larger and more embarrassing finding. Every application
imports the root package to obtain its dependency struct, so the root package's
import graph is the floor beneath all 107 of them. That floor was 1{,}798
packages, and a single import accounted for 1{,}154 of them.

The import was used for exactly four function-pointer registrations --- a tier
reader, a balance reader, a usage recorder, an ingest dialer --- all pointing
\emph{outward}. Nothing read a value back. A second import, a Kubernetes client
used by one file to poll the live pod set, accounted for 263 more.

Both inverted trivially once observed. The root package declares the callbacks
as plain function types and keeps the closures it builds; the composition root,
which already links the heavy dependencies, performs the registration. The pod
lister moved to its own package and is installed into an exported variable.

\begin{table}[h]
\centering
\begin{tabular}{lrr}
\toprule
\textbf{root package} & \textbf{before} & \textbf{after} \\
\midrule
total packages     & 1{,}798 & 644 \\
\texttt{k8s.io/}   & 263 & 0 \\
\texttt{aws/}      & 96  & 0 \\
\texttt{grpc}      & 65  & 0 \\
\texttt{go-git}    & 56  & 0 \\
\texttt{cloud.google.com/} & 18 & 0 \\
\bottomrule
\end{tabular}
\caption{Inverting two outward-only imports. Behaviour unchanged.}
\label{tab:floor}
\end{table}

The general lesson is worth stating plainly, because it is cheap to check and
expensive to miss: \emph{an import that only ever pushes values outward does not
need to be an import}. Declare the type locally, let the composition root do the
registration, and the dependency leaves the graph of everything downstream.

\section{Limits}

We report what these measurements do not establish.

The plugin figure of 646 packages still includes the 644-package root, because
the application takes the shared dependency struct. A plugin that took only what
it needs would be smaller again; we have not built one, and doing so requires
narrowing the dependency struct per application.

Eighty-three direct import edges exist \emph{between} application packages.
Each is a compile-time symbol reference that must become an interface, an RPC
call, or a merge before those two applications can occupy different processes.
Seven packages are hubs. This is the real remaining work and it is not
mechanical.

Several process-global couplings bind at load time rather than through the
dependency struct --- an indexer installed by one application into another, a
dispatcher assembled at the composition root to break an import cycle, a
single-process ledger resolved by every consumer of money. Each would break
silently, not loudly, across a process boundary.

The per-organisation database handle cache assumes one process per
\texttt{\{service\}.db} and is fenced by a single-writer lease. Applications
that share a service name cannot split without splitting that ownership first.

Finally, the numbers are warm-cache, single-machine, single-run. They measure
the developer loop, which is what we set out to improve; they are not
cold-build or CI-throughput measurements, and the ratios should not be quoted
as such.

\section{Conclusion}

The unified binary was the right call and remains so: one artifact, no mesh, no
sidecars, in-process calls between subsystems. What it accidentally also bought
was a shared link whose cost every developer pays on every edit, and a
dependency floor that one careless import raised for 107 applications at once.

Neither required abandoning the design. Making the composition primitive a
single function type meant a linked-in application and a separately built one
became the same thing, so the boundary could be moved per deployment without
touching code. Embedding the plugin binaries preserved the single artifact.
Inverting two outward-only imports cut the shared floor by 64\%.

The edit--rebuild loop went from 20.23\,s and 5.6\,GB to 2.11\,s and 0.63\,GB.
The interesting part is not the ratio; it is that the cost was never in the code
anyone was editing.

\end{document}
